\subsection{经期还能做什么：口交、肛交、自慰与"高潮能缓解痛经吗"}

经期性行为不等于插入式性交。以下三种方式在处理上差别很大，值得分开说明。

\textbf{一、口交}

\begin{itemize}
\item \textbf{一般没有禁忌}，主要障碍是观念上的不适。经血本身不是"毒"，也不会因为接触而对健康人构成伤害；
\item \textbf{但确实存在一条需要知晓的暴露路径}：经血是某些病原体的载体（乙肝、丙肝、HIV 等），若口腔一方存在破损——口腔溃疡、牙龈出血、近期拔牙或咽部手术——就构成了血液暴露。此时使用\textbf{口腔屏障膜（dental dam）}，或用剪开的避孕套、保鲜膜作为屏障，可以显著降低风险；
\item \textbf{顾虑较大时}，直接避开经量大的日子是完全可以接受的安排，不必勉强任何一方；
\item \textbf{不接受的一方直接说明即可}——不接受与不接受其他性行为一样，不需要理由。
\end{itemize}

\textbf{二、肛交}

\begin{itemize}
\item 经期做肛交并没有额外的解剖学禁忌，但\textbf{风险链更复杂}：同时存在经血暴露与肠道菌暴露，而且直肠黏膜本身容易在摩擦中损伤；
\item \textbf{顺序规则在经期尤其要严格}：肛交后的器具、手指或阴茎，未经更换屏障，\textbf{绝不能直接进入阴道}——经期宫颈口微开、内膜脱落，这条路径造成的上行感染风险比平时更高；
\item \textbf{全程使用安全套}，并在肛门与阴道之间更换；有痔、肛裂急性期、盆腔炎病史者应避免；
\item 若一方对"经血 + 肛交"的组合感到心理不适，这也是完全正当的拒绝理由。
\end{itemize}

\textbf{三、自慰}

\begin{itemize}
\item 经期自慰通常是最简单、最安全的选择：外阴刺激不增加感染风险，还可以完全由自己掌握节奏、环境与是否需要清理的顾虑；
\item \textbf{注意事前洗手}；如使用玩具，事后清洁或更换安全套；
\item 对经期容易感到不适、又不想完全放弃性愉悦的人来说，这往往是负担最小的方式——它不需要协商，也不需要"克服"什么。
\end{itemize}

\textbf{四、"高潮能缓解痛经吗"：机制、证据与边界}

这是一个流传很广的经验之谈，值得把它讲清楚：

\begin{itemize}
\item \textbf{可能机制}：高潮时释放的内啡肽与催产素有镇痛与放松作用；子宫在收缩后进入一段放松期；同时注意力转移、交感张力下降，也会让痛感的主观强度降低；
\item \textbf{证据强度}：有小规模研究与大量个人经验支持"高潮或自慰可短暂缓解痛经与痉挛"，但样本量小、个体差异大。用循证的语言说，它属于\textbf{"可能有效"}，而非已确证的疗法；
\item \textbf{边界一}：它不是治疗，不能替代就医。若痛经每月都严重影响上学、上班或日常生活，止痛药效果差，应排查子宫内膜异位症、子宫腺肌症、盆腔炎等继发性原因（参见本书常见月经问题相关小节）；
\item \textbf{边界二}：对一部分人来说，高潮后的子宫收缩会\textbf{短暂加重痉挛}。如果你发现自己属于这一类，就没有必要勉强用它来止痛——这不是"你没用对"，而是个体差异；
\item \textbf{边界三}：如果痛经伴随经量骤增、发热、非经期出血，这些都不是靠高潮能解决的问题，应当就诊。
\end{itemize}

\begin{tcolorbox}[colback=pink!5!white,colframe=red!50!black,title=贴心小叮咛]
把"经期能做什么"列清楚，是为了让人不被"经期什么都不能做"的暗示限制住；但列出选项不等于推荐。\textbf{知情之后的选择权，始终在你自己手里。}
\end{tcolorbox}

